% !TeX root = ../复分析笔记.tex
% 由主文档通过 \input 载入；请编译主文档。
\setcounter{chapter}{0} % 保留原书章号。
\chapter{绪论}
\label{chap:ca-introduction}

\section{研究对象与贯穿全书的问题}

本课程统一称为\emph{复分析（complex analysis）}，亦称复变函数论；本文档题名为《复分析笔记》。复分析研究复变量函数的微分、积分、局部展开与大域性质。本书以单复变量为范围，从复平面的几何出发，逐步解释一个核心现象：一个函数只要在开集上处处复可微，就受到远强于实可微性的约束。它在每一点附近都能展开为收敛幂级数，其内部取值与边界积分相联系，其零点、奇点以及映射形状也具有严格的结构。

全书围绕四组问题展开：复导数为什么要求不同方向的差商给出同一个极限？曲线积分何时与路径无关，何时能够恢复区域内部的函数？局部级数如何控制零点和奇点，并帮助计算实积分？局部定义的函数怎样延拓、逼近，或者把一个区域映到另一个区域？这些问题依次连接分析、平面拓扑与几何，而留数计算只是其中一条应用线索。

\section{从实变量走向复变量}

在实数范围内，方程 $x^2+1=0$ 没有解。引入满足 $i^2=-1$ 的虚数单位后，这个方程有了两个解 $i$ 与 $-i$，但复数的意义远不止于补齐代数方程的解。复数把平面上的两个实坐标组织成一个数，使伸缩、旋转、振荡等现象能够用统一的运算表示。研究复变量函数，也就是研究从复平面的一部分到复平面的映射，可以同时看到代数、几何与微积分之间的联系。

\emph{复变函数论（complex function theory）}的一个核心问题是：把实微积分中的自变量换成复数后，哪些结论仍然成立，哪些结论变得更强？例如，实函数 $x\mapsto x^2$ 的导数由沿实轴的差商决定，而复函数 $z\mapsto z^2$ 的导数要求从平面上任意方向趋近都得到同一个结果。这个要求看起来只改变了极限的取法，却会带来很强的约束：一个函数只要在一个开集内处处复可导，就自动有任意阶导数，并且在每一点附近等于自己的泰勒级数。实函数通常没有这样的性质。

另一个贯穿课程的问题是内部与边界的关系。对适当的复函数，知道一条闭曲线上的函数值，就能用积分恢复曲线内部的函数值；闭路积分往往只由少数奇点附近的一个系数决定。因此，复变函数既是一门研究函数结构的理论，也提供了计算实积分、判断多项式根的个数、求解平面边值问题的工具。学习时应把公式理解为结构的结果，尤其要辨认定义域、路径方向与奇点位置。

\section{先修知识、课程范围与阅读层次}

默认读者掌握一元微积分、实数列与函数项级数、多元函数的偏导数和全微分，以及二维向量和线性映射。开集、紧集、连通性、一致收敛和曲线积分将在需要时复习；抽象拓扑、测度论和泛函分析不作为核心部分的先修课程。证明中使用紧性、极限交换或积分估计时，均需说明具体依据。

本书按本科复分析课程的核心内容组织，并设置进阶选读部分。第~\ref{chap:ca-plane} 章至第~\ref{chap:ca-harmonic} 章构成首次学习的主体；第~\ref{chap:ca-normal} 章至第~\ref{chap:ca-approximation} 章进一步讨论存在性、延拓和构造问题。黎曼曲面只作入门，椭圆函数、模形式、多复变量与完整的值分布理论不在本书正文范围内。

本书将基本概念、核心证明与典型计算结合组织。正文优先说明结论为什么成立、条件在哪里使用；技术性推广另作提示。基础篇由复平面逐步进入积分理论、函数结构和共形几何，进阶篇完整展开正规族、黎曼映射、解析延拓、函数构造与全纯逼近。

\section{课程总纲}

下表给出全书的主题、核心问题和学习层次。“进阶”部分适合完成核心课程后选读。

\begin{longtable}{@{}p{0.07\textwidth}p{0.28\textwidth}p{0.49\textwidth}p{0.08\textwidth}@{}}
\toprule
章节 & 主题 & 核心问题 & 层次 \\
\midrule
\endhead
\ref{chap:ca-plane} & 复数与复平面 & 复数运算的几何意义；区域与曲线的基本语言 & 核心 \\
\ref{chap:ca-derivative} & 复可微性与全纯函数 & 复导数与实全微分的联系；柯西--黎曼条件 & 核心 \\
\ref{chap:ca-elementary} & 初等函数与分支 & 指数、对数、幂函数；多值性从何而来 & 核心 \\
\ref{chap:ca-integral} & 复积分与原函数 & 参数化、积分估计、路径无关与绕数 & 核心 \\
\ref{chap:ca-cauchy} & 柯西理论 & 积分定理、积分公式及其适用的区域条件 & 核心 \\
\ref{chap:ca-series} & 幂级数、零点与全纯函数结构 & 泰勒展开、恒等定理与最大模原理 & 核心 \\
\ref{chap:ca-laurent} & 洛朗级数与奇点 & 环域展开、奇点分类与无穷远点 & 核心 \\
\ref{chap:ca-residue} & 留数及积分应用 & 奇点如何决定闭路积分；实积分如何处理 & 核心 \\
\ref{chap:ca-counting} & 辐角原理与零点计数 & 不求出零点也能计数；扰动如何保持零点数 & 核心 \\
\ref{chap:ca-conformal} & 共形映射 & 局部保角与整体双全纯；典型区域的变换 & 核心 \\
\ref{chap:ca-harmonic} & 调和函数 & 调和共轭、泊松公式与圆盘边值问题 & 核心 \\
\ref{chap:ca-normal} & 正规族与黎曼映射 & 函数列的紧性如何产生共形映射 & 进阶 \\
\ref{chap:ca-continuation} & 解析延拓与黎曼曲面 & 局部函数如何延续；沿路径绕行会发生什么 & 进阶 \\
\ref{chap:ca-entire} & 整函数与亚纯函数构造 & 怎样指定零点或主部并构造函数 & 进阶 \\
\ref{chap:ca-approximation} & 全纯逼近与综合专题 & 怎样用有理函数逼近；如何综合使用全书工具 & 进阶 \\
\bottomrule
\end{longtable}

\section{知识依赖与学习路线}

基础主线是复数与拓扑语言、复可微性、曲线积分、柯西理论、幂级数、零点与奇点，最后进入留数和辐角原理。初等函数在主线前段提供可反复使用的例子；调和函数和共形映射在主线后段汇合几何与分析。

为避免证明上的循环依赖，第~\ref{chap:ca-derivative} 章先以复可微定义全纯性；第~\ref{chap:ca-elementary} 章借助实指数和实三角函数构造复指数；第~\ref{chap:ca-cauchy} 章从复可微性证明柯西积分理论；到第~\ref{chap:ca-series} 章再证明全纯函数具有局部幂级数展开。不能提前用“全纯函数已经可以展开为幂级数”证明柯西定理。

初次学习可顺序完成核心各章，并在每章检查三个层面：能否准确复述条件，能否解释证明中的关键步骤，能否处理一个不满足条件的反例。完成核心部分后，进阶各章按正文顺序学习；若主要关心函数构造，也可在第~\ref{chap:ca-continuation} 章之后直接集中学习第~\ref{chap:ca-entire} 章。

\section{记号约定}

以 $\mathbb{R}$、$\mathbb{C}$ 分别表示实数域、复数域，$i$ 表示满足 $i^2=-1$ 的虚数单位。通常写 $z=x+iy$，其中 $x,y\in\mathbb{R}$；$\overline z$、$|z|$、$\operatorname{Re}z$、$\operatorname{Im}z$ 分别表示共轭、模、实部和虚部。复函数通常写成 $f=u+iv$，其中 $u,v$ 是实值函数。

区域 $\Omega$ 指非空连通开集，$B(a,r)=\{z\in\mathbb{C}:|z-a|<r\}$ 表示以 $a\in\mathbb{C}$ 为中心、半径 $r>0$ 的开圆盘，$\mathbb{D}=B(0,1)$。$\partial\Omega$ 表示边界，$\overline\Omega$ 表示闭包，$K\Subset\Omega$ 表示 $K$ 是包含于 $\Omega$ 的紧集。曲线积分默认沿分段连续可微曲线进行；正向边界按行进方向使区域位于左侧，因此多连通区域的内边界取顺时针方向。

章末习题涵盖基本计算、证明与反例，基础篇的解答或提示列于附录~\ref{chap:ca-solutions}，进阶篇的解答随章给出。符号 $\log r$ 或 $\ln r$（$r>0$）表示实自然对数；复对数的分支另行定义。未经特别说明，函数等式均在两侧有定义的共同区域内讨论。

\section{参考学习材料}

课程覆盖范围可对照 MIT 的\href{https://ocw.mit.edu/courses/18-04-complex-variables-with-applications-spring-2018/pages/lecture-notes/}{18.04：Complex Variables with Applications（2018）}与\href{https://ocw.mit.edu/courses/18-112-functions-of-a-complex-variable-fall-2008/pages/lecture-notes/}{18.112：Functions of a Complex Variable（2008）}官方讲义目录。前者适合配合计算与应用，后者可用于核对积分理论、无穷乘积、正规族和黎曼映射等主题。进阶方向可参照 Harvard 的\href{https://people.math.harvard.edu/~ctm/home/text/class/harvard/213a/06/html/syl.html}{Math 213a（2006）课程大纲}。本书的章节顺序按上述知识依赖独立组织，不将这些课程的全部专题纳入正文。


